% 原书分章 PDF 第 8—11 页：第 22.3 节导言及第 22.3.1—22.3.2 节。
\section{数据自适应确定的动态处理下各样本切分特定反事实平均结局之平均的统计推断}

令 $\hat d:\mathcal M\to\mathcal D$ 为一个把经验分布映射为个体化处理规则的估计量。令 $B_n\in\{0,1\}^n$ 表示一次交叉验证切分的随机向量。对于切分 $B_n$，$P_{n,B_n}^0$ 表示训练样本 $\{i:B_n(i)=0\}$ 的经验分布，$P_{n,B_n}^1$ 表示验证样本 $\{i:B_n(i)=1\}$ 的经验分布。考虑 $J$ 折交叉验证方案。在 $J$ 折交叉验证中，数据以均匀随机方式划分成 $J$ 个互不相交且合起来覆盖全部数据的集合，每个集合的大小约为 $n/J$。每个集合依次用作一次验证集，其余所有集合的并集用作训练集。以概率 $1/J$，$B_n$ 在验证集 $j\in\{1,\ldots,J\}$ 的所有索引上取值为 1，在不对应于训练集 $j$ 的所有索引上取值为 0。

本节首先给出一个方法，为下列数据自适应目标参数提供估计量和统计推断：
\[
 \tilde\psi_{0n}=\E_{B_n}\Psi_{\hat d(P_{n,B_n}^0)}(P_0).
\]
注意，这个目标参数定义为若干数据自适应参数的平均；这些数据自适应参数从大小约为 $n(J-1)/J$ 的训练样本中学习得到。将估计量 $\hat d$ 分别应用于 $J$ 个训练样本中的每一个，就得到一个目标参数值 $\Psi_{\hat d(P_{n,B_n}^0)}(P_0)$；我们的目标参数 $\tilde\psi_{0n}$ 定义为这 $J$ 个目标参数的平均。

\subsection{CV-TMLE 的一般描述}

这里概述 CV-TMLE 的过程。在第 22.6 节中，我们将给出一个满足本节所述全部性质的具体 CV-TMLE。用 $j=1,\ldots,J$ 标记 $B_n$ 的各个实现值，并令 $d_{nj}=\hat d(P_{n,j}^0)$，其中 $\hat d$ 是最优规则的某个估计量。令
\[
 (a(0),\bar l(1))\mapsto
 \E_{nj}[Y\mid\bar A(1)=d_{nj}(a(0),v),\bar L(1)=\bar l(1)]
\]
表示基于训练样本 $j$ 对 $\E_0[Y\mid\bar A(1)=d_{nj}(A(0),V),\bar L(1)]$ 的初始估计。类似地，令
$l(0)\mapsto\E_{nj}[Y_{d_{nj}}\mid L(0)=l(0)]$
表示基于训练样本 $j$ 对 $\E_0[Y_{d_{nj}}\mid L(0)]$ 的初始估计。最后，令 $Q_{L(0),nj}$ 表示验证样本 $j$ 中 $L(0)$ 的经验分布。接着，使用下列子模型扰动这三个回归函数：
% 原书第三行集合末尾有两个可见右花括号，保持原样。
\[
\begin{gathered}
 \bigl\{\E_{nj}^{(\epsilon_2)}[Y\mid\bar A(1)=d_{nj}(a(0),v),
 \bar L(1)=\bar l(1)]:\epsilon_2\in\mathbb R\bigr\},\\
 \bigl\{\E_{nj}^{(\epsilon_1)}[Y_{d_{nj}}\mid L(0)=l(0)]:
 \epsilon_1\in\mathbb R\bigr\},\\
 \bigl\{Q_{L(0),nj}^{(\epsilon_0)}:\epsilon_0\in\mathbb R\}\bigr\},
\end{gathered}
\]
其中，这些子模型依赖基于训练样本 $j$ 得到的 $g_0$ 的估计 $g_{nj}$，并满足
\[
\begin{aligned}
 &\E_{nj}^{(0)}[Y\mid\bar A(1)=d_{nj}(a(0),v),\bar L(1)]\\
 &\qquad=\E_{nj}[Y\mid\bar A(1)=d_{nj}(a(0),v),\bar L(1)],\\
 &\E_{nj}^{(0)}[Y_{d_{nj}}\mid L(0)]
 =\E_{nj}[Y_{d_{nj}}\mid L(0)],\\
 &Q_{L(0),nj}^{(0)}=Q_{L(0),nj}.
\end{aligned}
\]
令 $Q_{nj}^{d_{nj}}(\epsilon)$ 表示给出上述三个经 $\epsilon\equiv(\epsilon_0,\epsilon_1,\epsilon_2)$ 扰动后的回归函数的参数映射。对固定的 $\epsilon$，$Q_{nj}^{d_{nj}}(\epsilon)$ 仅通过验证样本 $j$ 中 $L(0)$ 的经验分布依赖 $P_{nj}^1$。令 $\phi$ 为 $Q_0^d$ 的一个有效损失函数，使得 $Q_0^d=\argmin_{Q^d}P_0\phi(Q^d)$；并要求 $\phi$ 和上述子模型满足
\[
 D^*(d,Q^d,g)\in
 \left\langle\left.\frac{d}{d\epsilon}\phi(Q^d(\epsilon))\right|_{\epsilon=0}\right\rangle,
\]
其中，$\langle f\rangle=\{\sum_j\beta_jf_j:\beta\}$ 表示由 $f$ 的各分量张成的线性空间。选择 $\epsilon_n$，使其在 $\epsilon\in\mathbb R^3$ 上最小化 $P_n^1\phi(Q_{nj}^{d_{nj}}(\epsilon))$。然后定义 $Q_0^{d_{nj}}$ 的靶向估计
$Q_{nj}^{d_{nj}*}\equiv Q_{nj}^{d_{nj}}(\epsilon_n)$。
注意，在 M 估计量分析中通常采用的温和条件下，$Q_{nj}^{d_{nj}*}$ 保持 $Q_{nj}$ 的收敛率。我们需要 $\epsilon_n$ 及其对应更新 $Q_{nj}^{d_{nj}*}$ 具备的关键性质，是它（近似）解出有效影响曲线的交叉验证经验均值方程：
\begin{equation}
 \E_{B_n}P_{n,B_n}^1D^*(d_{nj},Q_{nj}^{d_{nj}*},g_{nj})
 =o_{P_0}(1/\sqrt n).
\tag{22.6}
\end{equation}
附录中给出的 CV-TMLE 实现满足这个方程，且可把 $o_{P_0}(1/\sqrt n)$ 替换为 0。所提出的 $\tilde\psi_{0n}$ 的估计量为
\[
 \tilde\psi_n^*\equiv\E_{B_n}\Psi_{d_{nj}}(Q_{nj}^{d_{nj}*}).
\]
第 22.6 节将给出计算 $\tilde\psi_n^*$ 的一个具体 CV-TMLE 算法；不过，对于不同的损失函数 $\phi$ 和子模型，也可以利用本节方法推导其他 CV-TMLE 算法。

\subsection{数据自适应参数\texorpdfstring{$\tilde\psi_{0n}$}{ψ̃₀ₙ}的统计推断}

现在分析 $\tilde\psi_{0n}$ 的 CV-TMLE $\tilde\psi_n^*$。首先给出 CV-TMLE 的一个表示定理。

\begin{bookresult}{定理}{22.4}
令 $g_{nj}$ 和 $d_{nj}$ 分别表示基于训练样本 $j$ 得到的 $g_0$ 和 $d_0$ 的估计。令 $Q_{nj}^{d_{nj}*}$ 表示第 22.3.1 节所给的 $Q_0^{d_{nj}}$ 的靶向估计，使得 $Q_{nj}^{d_{nj}*}$ 满足式 (22.6)。令 $R_{1d}$ 如定理 22.3 所定义。进一步假定，以趋于 1 的概率，$\max_jD^*(d_{nj},Q_{nj}^{d_{nj}*},g_{nj})$ 的上确界范数被某个 $M<\infty$ 所界定，并且
\[
 \max_{j\in\{1,\ldots,J\}}P_0
 \bigl\{D^*(d_{nj},Q_{nj}^{d_{nj}*},g_{nj})-D^*(d_1,Q^{d_1},g)\bigr\}^2
 \to0\quad\text{依概率},
\]
其中 $d_1\in\mathcal D$，而 $Q^{d_1}$ 和 $g$ 可能设定错误。最后，假定
\[
 \max_{j\in\{1,\ldots,J\}}
 \left|R_{1d_{nj}}(Q_{nj}^{d_{nj}*},Q_0^{d_{nj}},g_{nj},g_0)\right|
 =o_{P_0}(n^{-1/2}).
\]
则有
\[
 \tilde\psi_n^*-\tilde\psi_{0n}
 =(P_n-P_0)D^*(d_1,Q^{d_1},g^{d_1})+o_{P_0}(n^{-1/2}).
\]
\end{bookresult}

注意，上述定理中的 $d_1$ 不必等于最优规则 $d_0$；不过，稍后我们将讨论 $d_1=d_0$ 这个理想的特殊情形。上述定理也不要求 $g_0$ 已知，甚至不要求干预机制的极限 $g$ 等于 $g_0$。

还应注意，如果 $g_0$ 已知，使得所有 $g_{nj}$ 都可以被正确设定，那么立即可得
\[
 \max_{j\in\{1,\ldots,J\}}
 \left|R_{1d_{nj}}(Q_{nj}^{d_{nj}*},Q_0^{d_{nj}},g_{nj},g_0)\right|=0.
\]
实践中，即使 $g_0$ 已知，我们仍建议依据一个正确设定的模型估计它，因为这样可以提高效率（见 van der Laan 和 Robins，2003，第 2.3.7 节）。

如果上述定理的条件成立，渐近线性结果意味着
\[
 \sqrt n\bigl[\tilde\psi_n^*-\tilde\psi_{0n}\bigr]
 \to\operatorname{Normal}(0,\sigma_0^2),
\]
其中 $\sigma_0^2=P_0D^*(d_1,Q^{d_1},g^{d_1})^2$。在温和条件下，
\[
 \sigma_n^2=\frac1J\sum_{j=1}^JP_{n,j}^1
 \bigl\{D^*(d_{nj},Q_{nj}^{d_{nj}*},g_{nj})\bigr\}^2
\]
一致估计 $\sigma_0^2$。在 $\sigma_n^2$ 一致以及定理 22.4 的条件下，$\tilde\psi_{0n}$ 的一个渐近有效的 95\% 置信区间为
\begin{equation}
 \left[\tilde\psi_n^*\pm\frac{\sigma_n}{\sqrt n}\right].
\tag{22.7}
\end{equation}
